\documentclass[xcolor=table]{beamer}
\usepackage{beamerthemesplit}
\usepackage{wrapfig}
\usetheme{SPbGU}
\usepackage{pdfpages}
\usepackage{amsmath}
\usepackage{cmap} 
\usepackage[T2A]{fontenc} 
\usepackage[utf8]{inputenc}
\usepackage[english]{babel}
\usepackage{indentfirst}
\usepackage{amsmath}
\usepackage{tikz}
\usepackage{multirow}
\usepackage[noend]{algpseudocode}
\usepackage{algorithm}
\usepackage{algorithmicx}
\usepackage{fancyvrb}
\usetikzlibrary{shapes,arrows}
\usepackage{tikz}
\usetikzlibrary{automata,positioning}
\usepackage{listings}
\usepackage{multicol}
\setbeamertemplate{itemize items}[circle]

\lstdefinelanguage{scala}{
  morekeywords={%
          abstract,case,catch,class,def,do,else,extends,%
          false,final,finally,for,forSome,if,implicit,import,lazy,%
          match,new,null,object,override,package,private,protected,%
          return,sealed,super,this,throw,trait,true,try,type,%
          val,var,while,with,yield},
  otherkeywords={=>,<-,<\%,<:,>:,\#,@},
  sensitive=true,
  morecomment=[l]{//},
  morecomment=[n]{/*}{*/},
  morestring=[b]",
  morestring=[b]',
  morestring=[b]"""
}[keywords,comments,strings]

\lstset{language=Scala}

\newtheorem{mytheorem}{Theorem}
\renewcommand{\thealgorithm}{}


\beamertemplatenavigationsymbolsempty

\title[Parser Combinators for CFPQ]{Parser Combinators for Context-Free Path Querying}
% То, что в квадратных скобках, отображается в левом нижнем углу. 
\institute[JetBrains Research]{
JetBrains Research, Programming Languages and Tools Lab  \\
Saint Petersburg University
}

% То, что в квадратных скобках, отображается в левом нижнем углу.
\author[Kate Verbitskaia]{\textbf{Kate Verbitskaia}, Ilya Kirillov, Ilya Nozkin, Semyon Grigorev}

\date{September 28, 2018}

\definecolor{orange}{RGB}{179,36,31}

\begin{document}
{
\begin{frame}[fragile]
  \begin{tabular}{p{3.5cm} p{5.5cm} p{1cm}}
   \begin{center}
      \includegraphics[height=1.5cm]{pictures/jetbrainsResearch.pdf}
    \end{center}
    &
    \begin{center}
      \includegraphics[height=1.5cm]{pictures/icfplogo.png}
    \end{center}
    &
    \begin{center}
      \includegraphics[height=1.5cm]{pictures/SPbGU_Logo.png}
    \end{center} 
  \end{tabular}
  \titlepage
\end{frame}
}

\begin{frame}[fragile]
  \transwipe[direction=90]
  \frametitle{Path Querying}
\begin{minipage}[m]{0.42\linewidth}
\begin{tikzpicture}[shorten >=1pt,node distance=2cm,on grid,auto]
   \node[state] (a) {Here};
   \node[state] (b) [left=of a] {B};
   \node[state] (c) [right=of a] {C};
   \node[state] (d) [below=of c] {D};
   \node[state] (e) [below=of b] {E};
   \node[state] (f) [below=of a] {F};
   \node[state] (g) [below=of f] {There};
   \node[state] (h) [below=of d] {H};
    \path[->]
    (a) edge node {} (b)
    (a) edge node {} (c)
    (a) edge node {} (f)
    (b) edge node {} (e)
    (c) edge node {} (d)
    (d) edge node {} (h)
    (d) edge node {} (a)
    (e) edge node {} (f)
    (e) edge node {} (g)
    (f) edge node {} (b)
    (f) edge node {} (g)
    (h) edge node {} (g);    
\end{tikzpicture}
\end{minipage}\hfill
\begin{minipage}[m]{0.5\linewidth}
\pause    Is there a path from \textbf{Here} to \textbf{There}?
\end{minipage}
\end{frame}

\begin{frame}[noframenumbering]
  \transwipe[direction=90]
  \frametitle{Path Querying}
\begin{minipage}[m]{0.42\linewidth}
\begin{tikzpicture}[shorten >=1pt,node distance=2cm,on grid,auto]
   \node[state] (a) {Here};
   \node[state] (b) [left=of a] {B};
   \node[state] (c) [right=of a] {C};
   \node[state] (d) [below=of c] {D};
   \node[state] (e) [below=of b] {E};
   \node[state] (f) [below=of a] {F};
   \node[state] (g) [below=of f] {There};
   \node[state] (h) [below=of d] {H};
    \path[->]
    (a) edge[red] node {} (b)
    (a) edge node {} (c)
    (a) edge node {} (f)
    (b) edge[red] node {} (e)
    (c) edge node {} (d)
    (d) edge node {} (h)
    (d) edge node {} (a)
    (e) edge[red] node {} (f)
    (e) edge node {} (g)
    (f) edge node {} (b)
    (f) edge[red] node {} (g)
    (h) edge node {} (g);    
\end{tikzpicture}
\end{minipage}\hfill
\begin{minipage}[m]{0.5\linewidth}
    Is there a path from \textbf{Here} to \textbf{There}?
\end{minipage}
\end{frame}

\begin{frame}[noframenumbering]
  \transwipe[direction=90]
  \frametitle{Path Querying}
\begin{minipage}[m]{0.42\linewidth}
\begin{tikzpicture}[shorten >=1pt,node distance=2cm,on grid,auto]
   \node[state] (a) {Here};
   \node[state] (b) [left=of a] {B};
   \node[state] (c) [right=of a] {C};
   \node[state] (d) [below=of c] {D};
   \node[state] (e) [below=of b] {E};
   \node[state] (f) [below=of a] {F};
   \node[state] (g) [below=of f] {There};
   \node[state] (h) [below=of d] {H};
    \path[->]
    (a) edge node {} (b)
    (a) edge node {} (c)
    (a) edge node {} (f)
    (b) edge node {} (e)
    (c) edge node {} (d)
    (d) edge node {} (h)
    (d) edge node {} (a)
    (e) edge node {} (f)
    (e) edge node {} (g)
    (f) edge node {} (b)
    (f) edge node {} (g)
    (h) edge node {} (g);    
\end{tikzpicture}
\end{minipage}\hfill
\begin{minipage}[m]{0.5\linewidth}
    Is there a path from \textbf{Here} to \textbf{There} of length 13?
\end{minipage}
\end{frame}

\begin{frame}[noframenumbering]
  \transwipe[direction=90]
  \frametitle{Path Querying}
\begin{minipage}[m]{0.42\linewidth}
\begin{tikzpicture}[shorten >=1pt,node distance=2cm,on grid,auto]
   \node[state] (a) {Here};
   \node[state] (b) [left=of a] {B};
   \node[state] (c) [right=of a] {C};
   \node[state] (d) [below=of c] {D};
   \node[state] (e) [below=of b] {E};
   \node[state] (f) [below=of a] {F};
   \node[state] (g) [below=of f] {There};
   \node[state] (h) [below=of d] {H};
    \path[->]
    (a) edge node {} (b)
    (a) edge[red] node {} (c)
    (a) edge[red] node {} (f)
    (b) edge[red] node {} (e)
    (c) edge[red] node {} (d)
    (d) edge node {} (h)
    (d) edge[red] node {} (a)
    (e) edge node {} (f)
    (e) edge[red] node {} (g)
    (f) edge[red] node {} (b)
    (f) edge node {} (g)
    (h) edge node {} (g);    
\end{tikzpicture}
\end{minipage}\hfill
\begin{minipage}[m]{0.5\linewidth}
Is there a path from \textbf{Here} to \textbf{There} of length 13?
\end{minipage}
\end{frame}

\begin{frame}[noframenumbering]
  \transwipe[direction=90]
  \frametitle{Path Querying}
\begin{minipage}[m]{0.42\linewidth}
\begin{tikzpicture}[shorten >=1pt,node distance=2cm,on grid,auto]
   \node[state] (a) {Here};
   \node[state] (b) [left=of a] {B};
   \node[state] (c) [right=of a] {C};
   \node[state] (d) [below=of c] {D};
   \node[state] (e) [below=of b] {E};
   \node[state] (f) [below=of a] {F};
   \node[state] (g) [below=of f] {There};
   \node[state] (h) [below=of d] {H};
    \path[->]
    (a) edge node {} (b)
    (a) edge node {} (c)
    (a) edge node {} (f)
    (b) edge node {} (e)
    (c) edge node {} (d)
    (d) edge node {} (h)
    (d) edge node {} (a)
    (e) edge node {} (f)
    (e) edge node {} (g)
    (f) edge node {} (b)
    (f) edge node {} (g)
    (h) edge node {} (g);    
\end{tikzpicture}
\end{minipage}\hfill
\begin{minipage}[m]{0.5\linewidth}
    What is the shortest path from \textbf{Here} to \textbf{There}?
\end{minipage}
\end{frame}

\begin{frame}[noframenumbering]
  \transwipe[direction=90]
  \frametitle{Path Querying}
\begin{minipage}[m]{0.42\linewidth}
\begin{tikzpicture}[shorten >=1pt,node distance=2cm,on grid,auto]
   \node[state] (a) {Here};
   \node[state] (b) [left=of a] {B};
   \node[state] (c) [right=of a] {C};
   \node[state] (d) [below=of c] {D};
   \node[state] (e) [below=of b] {E};
   \node[state] (f) [below=of a] {F};
   \node[state] (g) [below=of f] {There};
   \node[state] (h) [below=of d] {H};
    \path[->]
    (a) edge node {} (b)
    (a) edge node {} (c)
    (a) edge[red] node {} (f)
    (b) edge node {} (e)
    (c) edge node {} (d)
    (d) edge node {} (h)
    (d) edge node {} (a)
    (e) edge node {} (f)
    (e) edge node {} (g)
    (f) edge node {} (b)
    (f) edge[red] node {} (g)
    (h) edge node {} (g);    
\end{tikzpicture}
\end{minipage}\hfill
\begin{minipage}[m]{0.5\linewidth}
    What is the shortest path from \textbf{Here} to \textbf{There}?
\end{minipage}
\end{frame}

\begin{frame}[fragile]
  \transwipe[direction=90]
  \frametitle{Path Querying: Regular Constraints}
\begin{minipage}[m]{0.42\linewidth}
\begin{tikzpicture}[shorten >=1pt,node distance=2cm,on grid,auto]
   \node[state] (a) {Here};
   \node[state] (b) [left=of a] {B};
   \node[state] (c) [right=of a] {C};
   \node[state] (d) [below=of c] {D};
   \node[state] (e) [below=of b] {E};
   \node[state] (f) [below=of a] {F};
   \node[state] (g) [below=of f] {There};
   \node[state] (h) [below=of d] {H};
    \path[->]
    (a) edge node[above] {f} (b)
    (a) edge node {f} (c)
    (a) edge node {l} (f)
    (b) edge node[left] {t} (e)
    (c) edge node {o} (d)
    (d) edge node {o} (h)
    (d) edge node[above] {i} (a)
    (e) edge node {t} (f)
    (e) edge node[below] {w} (g)
    (f) edge node[above] {a} (b)
    (f) edge node {e} (g)
    (h) edge node {d} (g);    
\end{tikzpicture}
\end{minipage}\hfill
\begin{minipage}[m]{0.5\linewidth}
\pause    Is there a path from \textbf{Here} to \textbf{There} of form $f \, o^*(d \mid l)$?
    
\bigskip 

\textcolor{white}{Constraints are a regular language}
\end{minipage}
\end{frame}

\begin{frame}[noframenumbering]
  \transwipe[direction=90]
  \frametitle{Path Querying: Regular Constraints}
\begin{minipage}[m]{0.42\linewidth}
\begin{tikzpicture}[shorten >=1pt,node distance=2cm,on grid,auto]
   \node[state] (a) {Here};
   \node[state] (b) [left=of a] {B};
   \node[state] (c) [right=of a] {C};
   \node[state] (d) [below=of c] {D};
   \node[state] (e) [below=of b] {E};
   \node[state] (f) [below=of a] {F};
   \node[state] (g) [below=of f] {There};
   \node[state] (h) [below=of d] {H};
    \path[->]
    (a) edge node[above] {f} (b)
    (a) edge[red] node {f} (c)
    (a) edge node {l} (f)
    (b) edge node[left] {t} (e)
    (c) edge[red] node {o} (d)
    (d) edge[red] node {o} (h)
    (d) edge node[above] {i} (a)
    (e) edge node {t} (f)
    (e) edge node[below] {w} (g)
    (f) edge node[above] {a} (b)
    (f) edge node {e} (g)
    (h) edge[red] node {d} (g);    
\end{tikzpicture}
\end{minipage}\hfill
\begin{minipage}[m]{0.5\linewidth}
    Is there a path from \textbf{Here} to \textbf{There} of form $f \, o^*(d \mid l)$?
    \pause
\bigskip 

Constraints are a regular language
\end{minipage}
\end{frame}

\begin{frame}[fragile]
  \transwipe[direction=90]
  \frametitle{Path Querying: Context-Free Constraints}
\begin{minipage}[m]{0.42\linewidth}
\begin{tikzpicture}[shorten >=1pt,node distance=2cm,on grid,auto]
   \node[state] (a) {Here};
   \node[state] (b) [left=of a] {B};
   \node[state] (c) [right=of a] {C};
   \node[state] (d) [below=of c] {D};
   \node[state] (e) [below=of b] {E};
   \node[state] (f) [below=of a] {F};
   \node[state] (g) [below=of f] {There};
   \node[state] (h) [below=of d] {H};
   \path[->]
    (a) edge node[above] {(} (b)
    (a) edge node        {)} (c)
    (a) edge node        {(} (f)
    (b) edge node[left]  {(} (e)
    (c) edge node        {(} (d)
    (d) edge node        {(} (h)
    (d) edge node[above] {)} (a)
    (e) edge node        {)} (f)
    (e) edge node[below] {)} (g)
    (f) edge node[above] {)} (b)
    (f) edge node        {)} (g)
    (h) edge node        {)} (g);    
\end{tikzpicture}
\end{minipage}\hfill
\begin{minipage}[m]{0.5\linewidth}
    Is there a path from \textbf{Here} to \textbf{There} which is a string of balanced brackets?
    
    \bigskip 
    
    Constraints are a context-free language given by the grammar:
    
    \begin{center}
    $ S \rightarrow \varepsilon \mid ( S ) S $
    \end{center}

\bigskip

\begin{center}
\textcolor{white}{(())()} \textcolor{white}{()} \textcolor{white}{()()}
\end{center}

\end{minipage}
\end{frame}

\begin{frame}[noframenumbering]
  \transwipe[direction=90]
  \frametitle{Path Querying: Context-Free Constraints}
\begin{minipage}[m]{0.42\linewidth}
\begin{tikzpicture}[shorten >=1pt,node distance=2cm,on grid,auto]
   \node[state] (a) {Here};
   \node[state] (b) [left=of a] {B};
   \node[state] (c) [right=of a] {C};
   \node[state] (d) [below=of c] {D};
   \node[state] (e) [below=of b] {E};
   \node[state] (f) [below=of a] {F};
   \node[state] (g) [below=of f] {There};
   \node[state] (h) [below=of d] {H};
   \path[->]
    (a) edge[red] node[above] {(} (b)
    (a) edge node        {)} (c)
    (a) edge[brown] node        {(} (f)
    (b) edge[red] node[left]  {(} (e)
    (c) edge node        {(} (d)
    (d) edge node        {(} (h)
    (d) edge node[above] {)} (a)
    (e) edge[red] node        {)} (f)
    (e) edge[red] node[below] {)} (g)
    (f) edge[red] node[above] {)} (b)
    (f) edge[brown] node        {)} (g)
    (h) edge node        {)} (g);    
\end{tikzpicture}
\end{minipage}\hfill
\begin{minipage}[m]{0.5\linewidth}
    Is there a path from \textbf{Here} to \textbf{There} which is a string of balanced brackets?
    
    \bigskip 
    
    Constraints are a context-free language given by the grammar:
    
    \begin{center}
    $ S \rightarrow \varepsilon \mid ( S ) S $
    \end{center}

\bigskip

\begin{center}
\textcolor{red}{(())()} \hspace{1cm} \textcolor{brown}{()} \hspace{1cm} ()()
\end{center}

\end{minipage}
\end{frame}


\begin{frame}[fragile]
  \transwipe[direction=90]
  \frametitle{Context-Free Path Querying}
  Given: 
  
  \begin{itemize}
    \item an input graph
    \item a context-free language specification
  \end{itemize}

\medskip
  
  Construct the set of paths in the graph which are in the language
\end{frame}

\begin{frame}[fragile]
  \transwipe[direction=90]
  \frametitle{Applications of Context-Free Path Querying}
\begin{itemize}
	\item Querying graph databases
	\item Static code analysis
	\item Network analysis
\end{itemize}
\end{frame}


\begin{frame}[fragile]
  \transwipe[direction=90]
  \frametitle{Parsing for Path Querying}

We need to parse strings along the paths

\bigskip 

Can we generalize parsing to work on graphs instead of strings?

\pause

\bigskip

Yes, we can!

\end{frame}

%
%\begin{frame}[fragile]
%  \transwipe[direction=90]
%  \frametitle{How to formulate the constraints?}
%
%As a string-based DSL: pass a string with a grammar to a dedicated processor
%
%\medskip
%  
%  \verb!graph.runQuery("S -> '(' S ')' S | "" ")!
%
%\medskip 
%
%Unsafe, not reusable, requires a step of parser generation.    
%\end{frame}
%
%\begin{frame}[fragile]
%  \transwipe[direction=90]
%  \frametitle{How to formulate the constraints?}
%
%As an ORM query
%
%\medskip
%  
%  \verb!graph.select(NAME)!
%
%\medskip 
%
%Tightly bound to a data model, not flexible.  
%\end{frame}
%
%\begin{frame}[fragile]
%  \transwipe[direction=90]
%  \frametitle{How to formulate the constraints?}
%
%With parser combinators!
%
%\medskip
%  
%  \verb!graph.runQuery("S -> '(' S ')' S | "" ")!
%
%\medskip 
%
%Cool! Safe! Highly reusable! wow!
%
%\end{frame}

 
\begin{frame}
  \transwipe[direction=90]
  \frametitle{Goals}

Writing queries should be user friendly!

It should be liberated from the particularities of parsing
 
\begin{itemize}
  \item Transparent integration of the query language
  \item Work with any context-free constraint
  \item Work with any input graph
\end{itemize}
  
\end{frame}

\begin{frame}[fragile]
  \transwipe[direction=90]
  \frametitle{Parser Combinators as a Way of Transparent Integration}
\begin{itemize}
  \item Encode the query in the programm itself
  \item Host language works for us: harder to write nonsensical grammar
  \item Composability 
  \item Natural representation of the result
\end{itemize}
\end{frame}

\begin{frame}[fragile]
  \transwipe[direction=90]
  \frametitle{Meerkat}
\begin{itemize}
  \item Generic parser combinator library in Scala
  \item CPS-parsing with memoization
  \item Supports any context-free grammar (left recursive, ambiguous)
  \item Semantics calculation
\end{itemize}

\begin{lstlisting}
val E: Nonterminal
= syn ( E ~ "^" ~ E
      | "-" ~ E
      | E ~ "*" ~ E 
      | E ~ "/" ~ E 
      | E ~ "+" ~ E 
      | E ~ "-" ~ E
      | "(" ~ E ~ ")"
      | "[0-9]".r 
      )
\end{lstlisting}
\end{frame}

\begin{frame}[fragile]
  \transwipe[direction=90]
  \frametitle{CFPQ in Meerkat}
\begin{minipage}[m]{0.42\linewidth}
\begin{tikzpicture}[shorten >=1pt,node distance=2cm,on grid,auto]
   \node[state] (a) {Here};
   \node[state] (b) [left=of a] {B};
   \node[state] (c) [right=of a] {C};
   \node[state] (d) [below=of c] {D};
   \node[state] (e) [below=of b] {E};
   \node[state] (f) [below=of a] {F};
   \node[state] (g) [below=of f] {There};
   \node[state] (h) [below=of d] {H};
   \path[->]
    (a) edge[red] node[above] {(} (b)
    (a) edge node        {)} (c)
    (a) edge[brown] node        {(} (f)
    (b) edge[red] node[left]  {(} (e)
    (c) edge node        {(} (d)
    (d) edge node        {(} (h)
    (d) edge node[above] {)} (a)
    (e) edge[red] node        {)} (f)
    (e) edge[red] node[below] {)} (g)
    (f) edge[red] node[above] {)} (b)
    (f) edge[brown] node        {)} (g)
    (h) edge node        {)} (g);    
\end{tikzpicture}
\end{minipage}\hfill
\begin{minipage}[m]{0.5\linewidth}
    Is there a path from \textbf{Here} to \textbf{There} which is a string of balanced brackets?
    
    \bigskip 
    
\begin{lstlisting}
val S: Nonterminal
= syn ( "(" ~ S ~ ")" ~ S
      | ""
      )
\end{lstlisting}

\bigskip

\begin{center}
\textcolor{red}{(())()} \hspace{1cm} \textcolor{brown}{()} \hspace{1cm} ()()
\end{center}

\end{minipage}
\end{frame}


\begin{frame}[fragile]
  \transwipe[direction=90]
  \frametitle{Supported Combinators}

\begin{table}[h]
\small
\centering
\begin{tabular}{l@{}|l}
\multicolumn{1}{c|}{Combinator} & \multicolumn{1}{|c}{Description} \\ \hline
{\lstinline!a ~ b!} & sequential parsing: {\lstinline!a!} then {\lstinline!b!}   \\
{\lstinline!a | b!} & choice: {\lstinline!a!} or {\lstinline!b!}         \\ \pause
{\lstinline!a ?!}   & optional parsing: {\lstinline!a!} or nothing   \\
{\lstinline!a *!}   & repetition of zero or more {\lstinline!a!} \\
{\lstinline!a +!}   & repetition of at least one {\lstinline!a!} \\ \hline \pause 
{\lstinline!a ^ f!} & apply {\lstinline!f!} function to {\lstinline!a!} if  {\lstinline!a!} is a token \\
{\lstinline!a ^^!}  & capture output of {\lstinline!a!} if {\lstinline!a!} is a token    \\
{\lstinline!a & f!} & apply {\lstinline!f!} function to {\lstinline!a!} if  {\lstinline!a!} is a parser \\
{\lstinline!a &&!}  & capture output of {\lstinline!a!} if {\lstinline!a!} is a parser    \\
\hline
\end{tabular}
\end{table}
\end{frame}

\begin{frame}[fragile]
  \transwipe[direction=90]
  \frametitle{How is it Done? Input}
\begin{center} String = linear graph \end{center} 
\begin{center}
\begin{tikzpicture}[shorten >=1pt,node distance=1.5cm,on grid,auto]
\tikzstyle{point} = [circle, minimum width = 3pt, fill]
   \node[point] (a) {};
   \node[point] (b)[right=of a] {};
   \node[point] (c)[right=of b] {};
   \node[point] (d)[right=of c] {};
   \node[point] (e)[right=of d] {};
   \node[point] (f)[right=of e] {};
   \node[point] (g)[right=of f] {};
   \path[->]
    (a) edge node {S} (b)
    (b) edge node {T} (c)
    (c) edge node {R} (d)
    (d) edge node {I} (e)
    (e) edge node {N} (f)
    (f) edge node {G} (g);
\end{tikzpicture}
\end{center}

\bigskip 

\bigskip 

\pause

\begin{center} What makes a graph \textbf{not} a string? \end{center}

\pause 

\begin{multicols}{2}
\begin{center} Branches \end{center}

\begin{multicols}{2}
\begin{center}
\begin{tikzpicture}[shorten >=1pt,node distance=1.5cm,on grid,auto]
\tikzstyle{point} = [circle, minimum width = 3pt, fill]
   \node[point] (a) {};
   \node[point] (b)[right=of a, above right=of a] {};
   \node[point] (c)[right=of a, below=of b] {};
   \path[->]
    (a) edge node {A} (b)
    (a) edge[below] node {B} (c);
\end{tikzpicture}
\end{center}

\columnbreak

\begin{center}
\begin{tikzpicture}[shorten >=1pt,node distance=1.5cm,on grid,auto]
\tikzstyle{point} = [circle, minimum width = 3pt, fill]
   \node[point] (a) {};
   \node[point] (b)[left=of a, above left=of a] {};
   \node[point] (c)[left=of a, below=of b] {};
   \path[->]
    (b) edge node {A} (a)
    (c) edge[below] node {B} (a);
\end{tikzpicture}
\end{center}
\end{multicols}

\columnbreak 

\pause

\begin{center} Cycles \end{center}
\begin{center}
\begin{tikzpicture}[shorten >=1pt,node distance=1.5cm,on grid,auto]
\tikzstyle{point} = [circle, minimum width = 3pt, fill]
   \node[point] (a) {};
   \node[point] (b)[right=of a] {};
   \node[point] (c)[below=of b] {};
   \node[point] (d)[left=of c] {};
   \path[->]
    (a) edge node {Y} (b)
    (b) edge node {A} (c)
    (c) edge node {D} (d)
    (d) edge node {A} (a);
\end{tikzpicture}
\end{center}
\end{multicols}
\end{frame}

\begin{frame}[fragile]
  \transwipe[direction=90]
  \frametitle{How is it Done? General Idea}
\begin{itemize}
  \item Each input position corresponds to a set of intermidiate parsing results
  \item Continue parsing along each branch independently
  \item Merge the sets of intermidiate results whenever two paths lead to the same input position
  \item Continue parsing only from the \textbf{new} results
  \item Memoize the results obtained 
\end{itemize}
\end{frame}


\begin{frame}[fragile]
  \transwipe[direction=90]
  \frametitle{Peeking into Vertices}
\begin{minipage}[m]{0.42\linewidth}
\begin{tikzpicture}[shorten >=1pt,node distance=2cm,on grid,auto]
   \node[state] (a) {Here};
   \node[state] (b) [left=of a] {B};
   \node[state] (c) [right=of a] {C};
   \node[state] (d) [below=of c] {D};
   \node[state] (e) [below=of b] {E};
   \node[state] (f) [below=of a] {F};
   \node[state] (g) [below=of f] {There};
   \node[state] (h) [below=of d] {H};
   \path[->]
    (a) edge[red] node[above] {(} (b)
    (a) edge node        {)} (c)
    (a) edge[brown] node        {(} (f)
    (b) edge[red] node[left]  {(} (e)
    (c) edge node        {(} (d)
    (d) edge node        {(} (h)
    (d) edge node[above] {)} (a)
    (e) edge[red] node        {)} (f)
    (e) edge[red] node[below] {)} (g)
    (f) edge[red] node[above] {)} (b)
    (f) edge[brown] node        {)} (g)
    (h) edge node        {)} (g);    
\end{tikzpicture}
\end{minipage}\hfill
\begin{minipage}[m]{0.5\linewidth}
    Is there a path from \textbf{Here} to \textbf{There} which is a string of balanced brackets?
    
    \bigskip 
    
\begin{lstlisting}
val Query : Nonterminal
= syn (LV("Here") ~ S ~ 
       LV("There"))

val S: Nonterminal
= syn ( "(" ~ S ~ ")" ~ S
      | ""
      )
\end{lstlisting}

\bigskip

\begin{center}
\textcolor{red}{(())()} \hspace{1cm} \textcolor{brown}{()} \hspace{1cm} ()()
\end{center}

\end{minipage}
\end{frame}

\begin{frame}[fragile]
  \transwipe[direction=90]
  \frametitle{Input abstraction}
\begin{lstlisting}
trait Input[+E, +N] {
  def checkNode(nodeId: Int,
      predicate: N => Boolean): Option[N]
  def filterEdges(nodeId: Int,
      predicate: E => Boolean): Seq[(E, Int)] }
\end{lstlisting}

\bigskip

We provide the integration for scalax.collection.Graph and Neo4j
\end{frame}

\begin{frame}[fragile]
  \transwipe[direction=90]
  \frametitle{Example}
\begin{lstlisting}
val query =
  syn((
    syn(LV("Actor") ^^) ~ outLE("ACTS_IN") ~
    LV("Movie")) & (a => (a.name, a.toInt)))
executeQuery(query, input)
  .groupBy{ case(a, i) => i }
  .toIndexedSeq
  .map{ case(i, ms) => (ms.head._1, ms.length) }
  .sortBy{ case(a, mc) => -mc }}
  .take (10)
\end{lstlisting}
\end{frame}

          
\begin{frame}
  \transwipe[direction=90]
  \frametitle{Evaluation}         
\begin{itemize}
  \item Ontology querying 
  \begin{itemize} 
    \item About 3 times faster than GLL-based solution
  \end{itemize}
  \item Static code analysis: may-alias relations
  \begin{itemize} 
    \item Up to 5 times faster than Trails
  \end{itemize}
  \item Movies database queries 
  \begin{itemize}
    \item Up to 10 times slower than Neo4j + Cypher
    \item All queries are regular constraints
  \end{itemize}
\end{itemize}

\end{frame}     

\begin{frame}
  \transwipe[direction=90]
  \frametitle{Limitations}
\begin{itemize}
  \item Overhead for the regular constraints
  \item Not exactly clear how to compute arbitrary semantics for the paths
  \begin{itemize}
    \item Paths can be lazily extracted, but in what order?
    \item Is it possible to compute some semantics in case of cycles?
  \end{itemize}
\end{itemize}
\end{frame}


\begin{frame}
  \transwipe[direction=90]
  \frametitle{Future Work}         
\begin{itemize}
  \item Reduce overhead for regular queries and improve performance
  \item Apply the library for static analysis of code
  \item Go beyond context-free constraints (ex. conjunctive grammars)
\end{itemize}  
\end{frame}     
            
            
\begin{frame}
  \transwipe[direction=90]
  \frametitle{Summary}  
Parser combinators for CFPQ:

\begin{itemize}
  \item Transparent integration of the query language
  \item Any context-free constraints
  \item Any input graph
\end{itemize}       
\end{frame}           
            
\begin{frame}
\transwipe[direction=90]
\frametitle{Contacts}
\begin{itemize}
  \item Semyon Grigorev: \href{mailto:s.v.grigoriev@spbu.ru}{s.v.grigoriev@spbu.ru}
  \item Kate Verbitskaia: \href{mailto:kajigor@gmail.com}{kajigor@gmail.com}
\end{itemize}
\begin{itemize}
  \item Meerkat for CFPQ: \href{https://github.com/YaccConstructor/Meerkat}{https://github.com/YaccConstructor/Meerkat}
\end{itemize}
\end{frame}
\end{document}
